\documentclass[sigconf,nonacm]{acmart}
\usepackage{booktabs}
% Generated by `python -m gapstudy.paper_numbers`. Do not edit by hand.
\providecommand{\TBD}{\textcolor{red}{\textbf{[TBD]}}}
\newcommand{\agentDirectedOverall}{\TBD}
\newcommand{\agentFlesch}{28.5}
\newcommand{\agentOnlyAgents}{\TBD}
\newcommand{\agentOnlyClaude}{\TBD}
\newcommand{\agentOnlyCopilot}{\TBD}
\newcommand{\agentOnlyOverall}{\TBD}
\newcommand{\agentUnitMedianWords}{39}
\newcommand{\agentWords}{608}
\newcommand{\anyGapOverall}{\TBD}
\newcommand{\contribFlesch}{51.0}
\newcommand{\contribWords}{445}
\newcommand{\detectorAUC}{\TBD}
\newcommand{\detectorBestMatcher}{\TBD}
\newcommand{\detectorFOne}{\TBD}
\newcommand{\fleschAgents}{37.6}
\newcommand{\fleschClaude}{18.6}
\newcommand{\fleschContribP}{$p < 0.0001$}
\newcommand{\fleschCopilot}{28.1}
\newcommand{\fleschGapContrib}{22.4}
\newcommand{\fleschHarderShare}{81\%}
\newcommand{\fleschPairs}{189}
\newcommand{\flowAgentFileIsASymlink}{1{,}965}
\newcommand{\flowAgentFileNotInEnglish}{1{,}869}
\newcommand{\flowAgentFileOnlyLinksElsewhere}{1{,}925}
\newcommand{\flowAllAgentFilesInPackage}{2{,}303}
\newcommand{\flowArchived}{2{,}123}
\newcommand{\flowCopyOfAnotherRepoSAgentFile}{1{,}835}
\newcommand{\flowFork}{2{,}200}
\newcommand{\flowNoContributingAtSnapshot}{755}
\newcommand{\flowNoReadmeAtSnapshot}{1{,}824}
\newcommand{\flowRepoNoLongerExists}{2{,}238}
\newcommand{\flowSnapshotCommitNotFound}{2{,}048}
\newcommand{\kappaCoverage}{\TBD}
\newcommand{\kappaWeighted}{\TBD}
\newcommand{\nAgentUnits}{2{,}664}
\newcommand{\nCodable}{2{,}661}
\newcommand{\nCodedUnits}{\TBD}
\newcommand{\nContribUnits}{1{,}862}
\newcommand{\nDocsFiles}{525}
\newcommand{\nDocsRepos}{113}
\newcommand{\nDocsSubsample}{30}
\newcommand{\nEligible}{755}
\newcommand{\nEligibleAgents}{214}
\newcommand{\nEligibleClaude}{252}
\newcommand{\nEligibleCopilot}{289}
\newcommand{\nFullyCodedRepos}{93}
\newcommand{\nInitUnedited}{12}
\newcommand{\nKnowledgeUnits}{\TBD}
\newcommand{\nMainRepos}{141}
\newcommand{\nPackageFiles}{2{,}303}
\newcommand{\nPerAgent}{67}
\newcommand{\nPilotRepos}{20}
\newcommand{\nPilotUnits}{161}
\newcommand{\nQueued}{1{,}564}
\newcommand{\nReadmeUnits}{2{,}594}
\newcommand{\nReliabilityRepos}{40}
\newcommand{\nReliabilityUnits}{338}
\newcommand{\nSampleRepos}{201}
\newcommand{\nStubContributing}{13}
\newcommand{\nUnits}{7{,}120}
\newcommand{\readmeFlesch}{27.8}
\newcommand{\readmeWords}{682}
\newcommand{\sampleMedianStars}{445}
\newcommand{\sharedOverall}{\TBD}


% Results that need frozen labels print as a red [TBD] (see numbers.tex).
% Regenerate the numbers with: python -m gapstudy.paper_numbers

\begin{document}

\title{What We Tell the Agent but Not the Newcomer}
\subtitle{Mining Agent Context Files for Gaps in Human Onboarding Documentation}

\author{Himanshu Mishra}
\affiliation{\institution{University of British Columbia}\city{Vancouver}\country{Canada}}
\email{himishra@student.ubc.ca}

\begin{abstract}
Coding agents such as Claude Code, OpenAI Codex and GitHub Copilot read a repository-level
context file before they work. Maintainers use these files to state how a project is built and
tested, how its code is organised and which conventions to follow --- knowledge a human newcomer
needs just as much. We ask how much of this knowledge is given only to the agent. From the
\nPackageFiles{} agent context files of the Agent READMEs dataset we sample \nSampleRepos{}
repositories that also have a CONTRIBUTING file (\nPerAgent{} per agent), recover the exact commit
of each agent file, and compare its \nCodable{} sections with the README and CONTRIBUTING at that
commit. Manual coding of \nQueued{} sections shows that \agentOnlyOverall{} of a repository's
agent-file knowledge is absent from its human documentation. % TODO: one sentence on RQ2 after coding.
Agent files are also harder to read than CONTRIBUTING files (median Flesch \agentFlesch{} vs
\contribFlesch{}). An automatic gap detector reaches F1 \detectorFOne{} on held-out repositories.
We release the dataset, codebook and detector.
\end{abstract}

\maketitle

\section{Introduction}

Coding agents start every session by reading a context file --- \texttt{CLAUDE.md},
\texttt{AGENTS.md} or \path{.github/copilot-instructions.md} --- that tells them how the project
is built and tested, how its code is organised and which conventions to follow~\cite{chatlatanagulchai2026agentreadmes,mohsenimofidi2026context}.
That is exactly what a newcomer to the project needs to know. People joining a project struggle
to learn its structure and unwritten rules~\cite{dagenais2010landscape}, and poor documentation is
among the barriers newcomers report most often~\cite{steinmacher2015barriers}.

Writing for an agent resembles writing for a newcomer: assume no prior context and make every
instruction explicit. This can push maintainers to write down knowledge that otherwise stays
implicit. A \texttt{CLAUDE.md} may say ``run \texttt{make verify} before committing; never edit
\texttt{src/gen/}'' while the \texttt{CONTRIBUTING.md} says only ``PRs welcome''. Context files are
most useful precisely when they capture such project-specific practice~\cite{gloaguen2026agentsmd}.
If this pattern is common, projects onboard their agents better than their human contributors.

Prior work describes what context files contain~\cite{chatlatanagulchai2026agentreadmes,mohsenimofidi2026context}
and whether they help agents~\cite{gloaguen2026agentsmd}. We ask what they reveal about a project's
documentation for \emph{people}:

\begin{itemize}
\item \textbf{RQ1:} What knowledge captured in agent context files is absent from the same repository's README and CONTRIBUTING files?
\item \textbf{RQ2:} Which newcomer needs and barriers does this agent-only knowledge address?
\item \textbf{RQ3:} How reliably can agent-only knowledge be identified automatically?
\end{itemize}

\noindent Our contributions are (1)~a measurement of how much project knowledge is documented only
for agents, (2)~a mapping of that knowledge to established newcomer needs and barriers, and (3)~an
open-source gap detector with a labelled dataset.

\section{Background and Related Work}

\paragraph{Agent context files.} \texttt{CLAUDE.md} is loaded at the start of every Claude Code
session; \texttt{AGENTS.md} is a tool-neutral convention first used by OpenAI Codex; Copilot reads
\path{.github/copilot-instructions.md}. Chatlatanagulchai et al.\ analysed 2{,}303 such files from
1{,}925 repositories, derived a taxonomy of 16 instruction categories and found the files hard to
read (median Flesch 16.6 for \texttt{CLAUDE.md})~\cite{chatlatanagulchai2026agentreadmes}.
Mohsenimofidi et al.\ studied how open-source projects engineer this context~\cite{mohsenimofidi2026context},
and Gloaguen et al.\ found that context files help agents mostly when they carry project-specific
practice rather than generic overviews~\cite{gloaguen2026agentsmd}.

\paragraph{Human documentation and newcomers.} README files describe what a project does and how
to install and use it~\cite{prana2019readme}; CONTRIBUTING files tell contributors how to set up,
test and submit changes. Dagenais et al.\ describe the ``landscape'' a newcomer must learn --- product,
processes and practices, team, documentation and context --- and the obstacles in
it~\cite{dagenais2010landscape}. Steinmacher et al.\ catalogue barriers newcomers face, including
orientation, technical hurdles and documentation problems~\cite{steinmacher2015barriers,santos2024onboarding}.
We use both frameworks to interpret agent-only knowledge.

\section{Study Design}

\subsection{Data}

\paragraph{Source.} We start from the replication package of Chatlatanagulchai et
al.~\cite{chatlatanagulchai2026agentreadmes}: \nPackageFiles{} agent files (922 \texttt{CLAUDE.md},
694 \texttt{AGENTS.md}, 687 Copilot) with their content, repository metadata and per-file commit
history. The package records each file's \emph{blob} SHA rather than a commit SHA (also in its
public mirror, under the name \path{content_commit_sha}). We therefore recover each file's
\emph{snapshot commit}: the newest commit in the file's history at which the file has that blob,
first from the package's history and otherwise from the GitHub API.

\paragraph{Human documentation.} At the snapshot commit we locate the README and CONTRIBUTING the
way GitHub does (\texttt{.github/}, then the root, then \texttt{docs/}), so the three documents of
a repository are mutually consistent: the human documentation as it stood when this version of
the agent file was committed. All \nSampleRepos{}$\times$3 documents were verified against their
git blob SHAs.

\paragraph{Exclusions and sampling.} Table~\ref{tab:flow} lists the exclusions in the order applied.
Link-only files point elsewhere (e.g.\ \texttt{@AGENTS.md}); near-copies (TF-IDF cosine $\ge 0.9$
with a file in another repository) are mostly tool-installed rule packs and mirrored projects;
language is detected with \texttt{langdetect} as in~\cite{chatlatanagulchai2026agentreadmes}.
Requiring a CONTRIBUTING file is by far the largest cut. From the \nEligible{} eligible files we drew
\nPerAgent{} per agent (seeded, one file per repository), giving \nSampleRepos{} repositories. Equal
allocation gives the per-agent comparisons equal power; overall estimates are post-stratified to the
eligible frame. Unedited \texttt{/init} output (\nInitUnedited{} files) is kept and examined in a
sensitivity analysis rather than excluded, because its body is generated from the repository's own
code.

\begin{table}[t]
\caption{Sampling flow (agent files).}
\label{tab:flow}
\small
\begin{tabular}{lr}
\toprule
Step & Remaining \\
\midrule
All agent files in the package & \flowAllAgentFilesInPackage \\
$-$ repository deleted & \flowRepoNoLongerExists \\
$-$ fork & \flowFork \\
$-$ archived & \flowArchived \\
$-$ snapshot commit not found & \flowSnapshotCommitNotFound \\
$-$ agent file is a symlink & \flowAgentFileIsASymlink \\
$-$ agent file only links elsewhere & \flowAgentFileOnlyLinksElsewhere \\
$-$ not in English & \flowAgentFileNotInEnglish \\
$-$ copy of another repository's file & \flowCopyOfAnotherRepoSAgentFile \\
$-$ no README at snapshot & \flowNoReadmeAtSnapshot \\
$-$ no CONTRIBUTING at snapshot & \flowNoContributingAtSnapshot \\
\midrule
Sampled (\nPerAgent{} per agent) & \nSampleRepos \\
\bottomrule
\end{tabular}
\end{table}

\subsection{Knowledge Units}

Each document is split with a CommonMark parser into sections at H2 and H3, plus H1 when a
document uses several H1s as sections; a single leading H1 is the title. Text before the first
section is a preamble unit; documents without section headings fall back to top-level blocks.
Tables of contents are flagged and not coded. This yields \nUnits{} units (\nAgentUnits{} agent,
\nReadmeUnits{} README, \nContribUnits{} CONTRIBUTING); agent units have a median of
\agentUnitMedianWords{} words. A pilot on 15 repositories showed sections to be mostly single topics, so
the section is the coding unit; checklists that are only partly documented are labelled
\emph{partially covered}.

\subsection{Coding}

To bound the effort we code at most 10 randomly chosen agent units per repository
(\nQueued{} units; \nFullyCodedRepos{} repositories coded in full), weighting each unit by the inverse
of its repository's sampling fraction. Each unit receives:
(A)~its instruction categories, using the 16 categories of~\cite{chatlatanagulchai2026agentreadmes};
(C)~its coverage in the README and CONTRIBUTING --- \emph{agent-only}, \emph{partially covered},
\emph{shared}, or \emph{agent-directed} for units that only address the AI (persona, response
style) and carry no project knowledge; paraphrases count, links elsewhere do not; the covering
human sections are recorded as evidence; and, for agent-only and partial units,
(D1)~the newcomer need, as one of Dagenais et al.'s landscape features~\cite{dagenais2010landscape},
and (D2)~the barriers it would mitigate~\cite{steinmacher2015barriers,santos2024onboarding}.
Flags record contradictions between agent and human documentation (\emph{conflict}) and topics
for which the human documentation only links elsewhere.

Coding used a purpose-built tool that shows the agent unit beside the full README and CONTRIBUTING
with keyword search. It deliberately shows no similarity ranking, so the gold labels are not
anchored on any matcher evaluated in RQ3. An open-source LLM could pre-fill categories (A) only;
coverage was always judged by the author. The codebook was refined on \nPilotRepos{} pilot repositories
(\nPilotUnits{} units). Reliability is intra-rater: the \nReliabilityRepos{} repositories of the
reliability set (\nReliabilityUnits{} units) were re-coded blind at least 14 days after the first pass,
with first-pass labels hidden and the order shuffled. We report Cohen's $\kappa$ and linearly
weighted $\kappa$~\cite{cohen1968weighted} on the ordinal coverage scale: $\kappa = \kappaCoverage$,
$\kappa_w = \kappaWeighted$.

\subsection{Analysis}

\paragraph{RQ1.} For each repository we compute the share of its knowledge units (excluding
agent-directed) that are agent-only, and average over repositories per agent; the overall figure
is post-stratified to the eligible frame (\nEligibleClaude{} \texttt{CLAUDE.md}, \nEligibleAgents{}
\texttt{AGENTS.md}, \nEligibleCopilot{} Copilot files). 95\% confidence intervals come from 1{,}000
bootstrap draws that resample repositories within each agent. Because units are nested in
repositories, differences across categories and agents are tested with a logistic GEE clustered by
repository (exchangeable working correlation, robust standard errors)~\cite{liang1986gee}; $\chi^2$
tests are reported only descriptively.

\paragraph{RQ2.} We report the distribution of needs (D1) and barriers (D2) among agent-only and
partially covered units, with examples.

\paragraph{Sensitivity.} We recompute RQ1 without unedited \texttt{/init} files, without the
\nStubContributing{} stub CONTRIBUTING files ($<$50 words, typically a link to a website), and without
topics the human docs only link out for. \nDocsRepos{} of the \nSampleRepos{} repositories also keep
Markdown under \texttt{docs/}; for a seeded subsample of \nDocsSubsample{} repositories
(\nDocsFiles{} files) we estimate with the tuned detector how many agent-only units \texttt{docs/}
would cover.

\paragraph{RQ3: gap detector.} Units are split into items (list entries, sentences, code blocks), so
a checklist is matched rule by rule. Each agent unit is scored against its repository's README and
CONTRIBUTING by (i)~TF-IDF, (ii)~three sentence-embedding models~\cite{reimers2019sbert} ---
all-MiniLM-L6-v2, bge-base-en-v1.5~\cite{xiao2023cpack} and Qwen3-Embedding-0.6B~\cite{zhang2025qwen3embedding}
--- at both section and item level, and (iii)~embeddings plus an LLM judge
(Qwen3-Coder-30B-A3B-Instruct~\cite{yang2025qwen3}, served with vLLM~\cite{kwon2023vllm}) that
sees the unit and its five best-matching human sections and decides only for borderline scores. All
model revisions are pinned. Score variants, thresholds and the LLM band are tuned on the pilot
repositories only; we report precision, recall and F1 for agent-only, ROC-AUC and three-class
macro-F1 on the remaining repositories, with repository-bootstrap confidence intervals.

\section{Results}

\subsection{Agent files are harder to read than CONTRIBUTING files}

Across all \nSampleRepos{} repositories (Table~\ref{tab:read}), agent files are about as long and
as hard to read as READMEs, but harder to read than the same repository's CONTRIBUTING file in
\fleschHarderShare{} of repositories (median difference \fleschGapContrib{} Flesch points,
\fleschPairs{} pairs, Wilcoxon \fleschContribP). By agent, median Flesch scores are \fleschClaude{}
(\texttt{CLAUDE.md}), \fleschCopilot{} (Copilot) and \fleschAgents{} (\texttt{AGENTS.md}); the
\texttt{CLAUDE.md} figure is close to the 16.6 reported by~\cite{chatlatanagulchai2026agentreadmes}.
Readability is computed with \texttt{textstat}~\cite{flesch1948readability} on prose with code removed.

\begin{table}[t]
\caption{Length and readability per document (medians; higher Flesch = easier).}
\label{tab:read}
\small
\begin{tabular}{lrr}
\toprule
Document & Words & Flesch \\
\midrule
Agent file & \agentWords & \agentFlesch \\
README & \readmeWords & \readmeFlesch \\
CONTRIBUTING & \contribWords & \contribFlesch \\
\bottomrule
\end{tabular}
\end{table}

\subsection{RQ1: How much agent knowledge is agent-only?}

Of \nCodedUnits{} coded units, \nKnowledgeUnits{} carry project knowledge; \agentDirectedOverall{} of
units address only the agent. On average, \agentOnlyOverall{} of a repository's agent-file
knowledge is absent from its README and CONTRIBUTING (\texttt{CLAUDE.md} \agentOnlyClaude{},
\texttt{AGENTS.md} \agentOnlyAgents{}, Copilot \agentOnlyCopilot{}); counting partial coverage,
\anyGapOverall{} is not fully available to human readers, and \sharedOverall{} is fully shared.
% TODO after coding: per-category table (analysis/phase7/rq1_by_category.csv), GEE odds ratios
% (rq1_gee.csv), and the sensitivity table (sensitivity.csv). Report which categories are most
% often agent-only and whether agents differ.

\subsection{RQ2: Which newcomer needs does it address?}

% TODO after coding: needs table (rq2_needs.csv), barriers table (rq2_barriers.csv), and 3--4
% examples from analysis/phase7/results.md.
\TBD

\subsection{RQ3: Can gaps be detected automatically?}

The best matcher on the pilot repositories, \detectorBestMatcher{}, reaches F1 \detectorFOne{} for
agent-only units on held-out repositories (ROC-AUC \detectorAUC).
% TODO after coding: full table from analysis/phase6/detector_results.md, the effect of the LLM
% stage, evidence hit@k, and the main error types from analysis/phase6/errors.md.

\section{Discussion}

% TODO after results. Planned points: (1) whether projects onboard agents better than people, and
% for which kinds of knowledge; (2) agent files as a source for improving CONTRIBUTING, via gapcheck;
% (3) readability: agent-only knowledge is also written in a harder register.
\TBD

\section{Threats to Validity}

\paragraph{Construct.} Coverage is judged against README and CONTRIBUTING only. Knowledge in
\texttt{docs/}, wikis or project websites is invisible to it; more than half of the repositories
(\nDocsRepos{}/\nSampleRepos{}) keep Markdown under \texttt{docs/}, and the \texttt{docs/}
sensitivity analysis estimates the effect. Units that only address the agent are separated out so
that they do not inflate the agent-only rate.

\paragraph{Internal.} One person coded all units; reliability is intra-rater (blind re-code after 14
days), which is weaker evidence than agreement between coders. LLM category suggestions, if used,
may anchor the coder on categories, but never on coverage. Fork and archive status were checked at
collection time, not at the snapshot.

\paragraph{External.} Requiring a CONTRIBUTING file selects larger, more established projects: the
sample's median is \sampleMedianStars{} stars versus 51--96 in the source dataset. Results describe
English-language repositories with all three documents at one point in time; agent files and their
conventions are evolving quickly.

\paragraph{Conclusion.} Units are nested in repositories; we therefore report repository-level
estimates with repository bootstrap intervals and use clustered models instead of plain $\chi^2$
tests.

\section{Conclusion}

% TODO: summarise RQ1--RQ3 findings once available.
We compared what \nSampleRepos{} projects tell their coding agents with what their README and
CONTRIBUTING tell people. \TBD

\section*{Data Availability}
Code, codebook, coding tool, labels and the \texttt{gapcheck} detector:
\url{https://github.com/himans-iitk/CS507A} (archived release: DOI \TBD).

\section*{Use of AI Tools}
% EDIT: describe your own use accurately. The draft below reflects how the project was built.
The study's data pipeline, coding tool, analysis code and a first draft of this paper were written
with the help of an AI coding assistant (Claude); the author reviewed and is responsible for all
of it. All coverage, need and barrier labels were assigned by the author. An open-source LLM was
used for optional category suggestions and as one component of the evaluated gap detector.

\bibliographystyle{ACM-Reference-Format}
\bibliography{refs}

\end{document}
